% =============================================================
% Qualitative findings — main-text version (NeurIPS-tight).
% Each finding is multi-source, multi-hop, and not derivable from
% a single document. Long lists, snapshot-level details, exact mix
% weights, and additional examples are deferred to the appendix.
% =============================================================

\paragraph{Hidden external-model dependencies.}

A core advantage of recursive dependency tracing is that it surfaces upstream models that never appear in the final model's own paper or card. These dependencies are often mediated by intermediate datasets, filters, classifiers, teachers, or tools, and therefore only become visible after following several evidence-backed hops across documents.

\begin{itemize}
    \item \textbf{SmolLM3 inherits a Llama-3-70B-mediated pretraining dependency through FineMath.}
    SmolLM3 documents FineMath as part of its pretraining mixture, and the FineMath card describes the use of a classifier to score and filter mathematical web data. Only by following the classifier card does the upstream dependency become visible: the \texttt{finemath-classifier} was trained on annotations generated by \texttt{Llama-3-70B-Instruct}. Thus, SmolLM3's pretraining pipeline is transitively shaped by Llama-3 through an intermediate filter, even though the SmolLM3 card does not name Llama-3 as a data source.

    \item \textbf{DR-Tulu inherits a Claude dependency through ScholarQA trajectories.}
    The DR-Tulu paper states that Ai2 ScholarQA trajectory data was used to create SFT data, but does not name the model used to generate those trajectories. Following the ScholarQA implementation reveals that its generation pipeline uses Claude Sonnet 3.7 by default. The graph therefore surfaces a transitive model-mediated dependency: \texttt{Claude Sonnet 3.7} $\rightarrow$ \texttt{Ai2 ScholarQA trajectories} $\rightarrow$ \texttt{DR-Tulu SFT}.

    \item \textbf{Nemotron-3-Content-Safety inherits DeepSeek-R1-generated content through NemoGuard.}
    The Nemotron-3-Content-Safety card describes the model as extending Llama-3.1-NemoGuard 8B Content Safety. The NemoGuard card identifies its base/fine-tuning lineage, while the NaturalReasoning documentation indicates that a subset of training questions was answered using DeepSeek-R1. Combining these sources reveals a transitive DeepSeek-R1-generated-data dependency that is not visible from the Nemotron-3-Content-Safety card alone.

    \item \textbf{OLMo-3 RL-Zero variants inherit Qwen2.5-Coder through intermediate checkpoints and data construction.}
    OLMo-3 RL-Zero models are trained on Dolci RL-Zero mixtures, which derive from earlier OLMo-3 checkpoints and code-oriented data construction. Tracing those intermediate artifacts back to the OLMo-3 midtraining/data-generation pipeline surfaces Qwen2.5-Coder-32B-Instruct as a transformation model used upstream. The dependency is not stated in the downstream RL-Zero cards, but emerges by following checkpoint and dataset lineage.
\end{itemize}

These examples show why final-model documentation is insufficient for model-mediated dependencies. A model card can accurately describe the immediate dataset or teacher while omitting the upstream generator, judge, or filter that shaped that artifact several steps earlier. Additional cases, including cross-organization base-model chains and judge-bootstrap dependencies, are deferred to the appendix.

\paragraph{Train/eval coupling and decontamination.}

The graph also surfaces relationships between training data and evaluation benchmarks that are not necessarily direct test-set leakage. Instead, these are cases of structural train/eval coupling: benchmark training splits, auxiliary data, prompts, validation splits, or concepts are transformed into training data while the same benchmark family remains an evaluation target. Such relationships can evade surface-form decontamination because the relevant connection is lineage-based rather than text-overlap-based.

\begin{itemize}
    \item \textbf{IFEval prompts are used in OLMo-3 training and IFEval remains an evaluation target.}
    OLMo-3 uses IF-RLVR prompts sampled from IFEval and IFBench-Train, while IFEval is also used in evaluation. This creates a direct audit question about whether evaluation prompts were excluded from training prompts.

    \item \textbf{SWE-bench and SWE-rebench appear in both training and evaluation roles within the Nemotron family.}
    Nemotron training artifacts include SWE-bench/SWE-rebench-derived resources, while Nemotron model cards or evaluation code report evaluations on SWE-bench-style benchmarks. This is stronger than cross-model reuse because the train/eval roles occur within the same product family.

    \item \textbf{LiveCodeBench validation appears in training artifacts while LiveCodeBench remains an evaluation benchmark.}
    Nemotron training artifacts include a \texttt{livecodebench\_v5\_validation} resource. When paired with LiveCodeBench evaluation use in related releases, this surfaces a benchmark-split audit question that would be hard to see without aligning training YAMLs and eval configs.

    \item \textbf{Train/eval dual-role is systemic across popular benchmarks.}
    Aggregating roles per benchmark, GSM8K appears in 127 evaluation edges and 25 training edges, MMLU in 135 / 7, GPQA in 110 / 4, Hendrycks MATH in 69 / 18, IFEval in 66 / 11, and SWE-bench Verified in 31 / 7. The graph makes the train/eval coupling visible at scale, where it is otherwise spread across hundreds of independent cards.
\end{itemize}

These cases do not by themselves establish benchmark contamination. Rather, they show that evaluation independence is a graph property: to assess it, one must trace whether benchmark artifacts appear upstream as seeds, prompts, validation splits, concept sources, or synthetic-data inputs. The graph also surfaces explicit decontamination receipts as positive counterexamples, which we discuss in the community-practice paragraph below.

\paragraph{License/provenance obligations propagate through synthetic-data pipelines.}

License risk is not always visible from the license attached to a final dataset or model. Prior work argues that dataset legal risk cannot be assessed from license terms alone and instead requires tracing redistribution and lifecycle history~\cite{kim2025trustlicensesseedataset}; supply-chain audits similarly find that license labels can function as weak metadata signals when the underlying compliance payload is missing or fails to propagate~\cite{jewitt2026permissivewashingopenaisupply}. Our recovered graphs surface analogous risks in LLM pipelines, where obligations may originate from models used to generate, label, or filter a dataset rather than only from the final model card.

\begin{itemize}
    \item \textbf{License-bearing root models reach 150+ downstream artifacts each.}
    Counting unique downstream artifacts reachable from any model in a license family yields: Qwen 184 descendants, GPT-4 178, Llama 180, DeepSeek 167, Claude 61, Gemini 34. License obligations from any one of these families therefore touch a substantial fraction of every recovered training pipeline. The reach is not directly stated anywhere; it emerges only from cross-document traversal.

    \item \textbf{OLMo-3's use of Llama-Nemotron Post-Training data creates a Llama-license audit surface.}
    The OLMo-3 paper lists Llama-Nemotron Post-Training data as part of the training mixture, while the NVIDIA dataset card provides more precise generator and license information. Only by joining the OLMo mixture description with the upstream dataset card does the Llama Community License surface for this slice of the training pipeline.

    \item \textbf{PrismMath contributes to Nemotron-3-Super through a Qwen-license-bearing dataset.}
    Nemotron-3-Super training artifacts reference PrismMath, and the PrismMath card warns that it may be subject to redistribution and use requirements in the Qwen License Agreement. The graph therefore surfaces a Qwen-related license-obligation path that is not visible from the final Nemotron model card alone.

    \item \textbf{Nemotron pretraining data contains stacked generator-license surfaces.}
    Nemotron pretraining datasets such as \texttt{Nemotron-Pretraining-Specialized-v1}, \texttt{Nemotron-CC-v2}, and \texttt{Nemotron-CC-Math-v1} explicitly name obligations from Qwen, DeepSeek, and Phi-4. The graph makes these obligations queryable as a stacked license/provenance surface inherited by downstream Nemotron models.
\end{itemize}

We treat these as license-obligation findings rather than legal determinations. \system{} identifies where obligations may propagate.

\paragraph{Hidden judge, filter, and synthesizer stacks shape training data.}

Many dependencies do not appear as explicit training datasets. Instead, upstream models act as judges, filters, classifiers, reward models, embedding models, or synthetic-data generators. These roles are often compressed in papers but described only in scripts, dataset cards, or manifests.

\begin{itemize}
    \item \textbf{Qwen3-32B is a pervasive judge across OLMo-3 RL variants.}
    The OLMo-3 paper notes the use of Qwen3-32B as an LM judge, while training scripts show that this judge is used across multiple RL-Zero specializations and post-training variants. A graph query over judge-role edges shows that GPT-4.1 and Qwen3-32B each recur across 20+ distinct OLMo-3 stages --- a per-stage multiplicity that is only visible after aggregating across YAMLs, dataset cards, and paper sections.

    \item \textbf{OLMo-3 uses multiple Gemma-3 variants for pretraining filtering.}
    The OLMo-3 paper mentions Gemma-based document classification, but the data-pipeline scripts reveal multiple distinct Gemma variants across model sizes and instruction/pretraining forms. The graph therefore exposes the full filtering panel rather than a single family-level mention.

    \item \textbf{Nemotron-3-Super uses DeepSeek-V3.2 as a tool-call trajectory judge.}
    The main Nemotron-3 report names some synthetic-data and judge models, but the agentic dataset card identifies DeepSeek-V3.2 as a judge for filtering tool-call trajectories. The graph therefore surfaces an additional judge dependency hidden in dataset-level documentation.

    \item \textbf{Nemotron-3-Super's RL signal is an entirely Qwen-bootstrapped reward stack.}
    \texttt{nvidia/Qwen3-Nemotron-235B-A22B-GenRM-2603} --- the sole reward signal across Super's RLHF and multi-environment RLVR --- is initialized from Qwen3-235B-A22B-Thinking-2507 and fine-tuned on \texttt{nvidia/HelpSteer3} and commercially-friendly subsets of \texttt{lmarena-ai/arena-human-preference-140k}. The four-hop chain (Qwen3 base $\rightarrow$ GenRM judge $\rightarrow$ Nemotron-3-Super) requires the Nemotron paper, the GenRM card, and the underlying preference-data card to assemble.

    \item \textbf{The graph's top generator/judge hubs are dominated by Qwen, DeepSeek, OpenAI, Meta, and Mistral models.}
    Aggregating across all model-mediated relations, the most-used upstream models include Qwen2.5-32B-Instruct (68 uses), DeepSeek-R1 (51), Llama-3.3-70B-Instruct (49), Qwen3-32B (42), and GPT-4.1 (41). This leaderboard reframes ``open'' AI2 / HuggingFace releases as products that are deeply downstream of partly-closed model families, and shows that cross-organization synthetic-data construction concentrates at the post-training stages (DPO, RLHF, RLVR).
\end{itemize}

\paragraph{Released artifacts can differ from trained-on artifacts, and documentation drifts from code.}

Artifact names alone are not stable evidence of what was used in training. A graph with per-operation evidence makes it easier to identify when released artifacts differ from the training-time artifacts they are meant to represent.

\begin{itemize}
    \item \textbf{The OLMo-3-7B reproduction dataset was modified after training.}
    The dataset intended to reproduce OLMo-3-7B-1025 was modified after training: some olmOCR science PDFs were redacted and replaced with \texttt{[REMOVED]}. The dataset card warns that this affects reproducibility, but the model paper does not foreground the issue.

    \item \textbf{Nemotron RL-Training-Blends use a placeholder-resolution pattern that recurs across product lines.}
    Nemotron-Super-RL-Training-Blends includes placeholder rows for DAPO-Math and Skywork-OR1 contributions, requiring a separate \texttt{fill\_placeholders.py} script to fetch and restore upstream data before the blend is usable. Similar placeholder mechanisms appear in Nemotron-3-Nano and Nemotron-3-Nano-Omni, suggesting a recurring release strategy in the Nemotron family rather than a one-off artifact issue.

    \item \textbf{SmolLM3's pretraining mixture is much more granular than its card summary.}
    The SmolLM3 card summarizes pretraining with broad categories such as web, code, and math data. The YAMLs reveal many specific subsets, language-specific substitutions, and stage-specific mixture weights. The card is useful for high-level understanding, but the YAMLs are necessary for reproduction.

    \item \textbf{SmolLM3's MegaMath source includes a Qwen-generated subset hidden behind a coarse dataset name.}
    The card names MegaMath as a single source, while code-level configuration distinguishes web-derived MegaMath data from Qwen-generated QA pairs. This reveals an additional model-generated subset that is obscured by the coarse dataset label.

    \item \textbf{OLMo-3 release artifacts contain smaller but concrete metadata conflicts.}
    OLMo-3 artifacts disagree on details such as the number of checkpoints used for souping, model-card prose sometimes names the wrong sibling dataset, and per-source prompt counts differ across related dataset cards. These examples show that even highly open releases can contain local documentation inconsistencies.
\end{itemize}

\paragraph{The graph also surfaces deliberate community good practice.}

Not every multi-hop pattern in the recovered graph is a risk. Several findings instead show developers acting carefully across hops --- swapping out license-encumbered teachers, decontaminating against benchmarks they did not author, or routing reward signals through commercially-friendly subsets. These choices are themselves multi-document patterns that no single card explicitly highlights, and the graph is what surfaces them as deliberate practice rather than coincidence.

\begin{itemize}
    \item \textbf{AI2's CraneMath, CraneCode, and MegaMatt re-run upstream recipes with permissive teachers.}
    Three OLMo-3 midtraining subsets share the same design pattern: CraneMath is ``an open-licensed reproduction of the SwallowMath recipe (Fujii et al., 2025): same FineMath4+ source and SwallowMath rewrite prompt, but generated with Qwen3 instead of Llama for permissive licensing.'' CraneCode applies the same swap to SwallowCode's two-stage SGCR$+$SCOR rewriting pipeline using Qwen2.5-Coder-32B-Instruct in place of Llama. MegaMatt independently recreates OctoThinker's MegaMath-Web-Pro-Max methodology with Qwen3, again to obtain a permissively-licensed alternative. The repeated pattern --- replicate the recipe, swap the teacher, document the swap --- only emerges by joining the OLMo-3 paper, the OLMo-3 dataset cards, and the SwallowMath / SwallowCode / MegaMath upstream cards.

    \item \textbf{Nemotron-3's GenRM training was deliberately restricted to commercially-friendly preference subsets.}
    The Nemotron-3-Super GenRM card attributes its training to ``commercially friendly subsets of lmarena-140k (Chiang et al. 2024)'' rather than the full lmarena release. This narrow subset choice is documented at the GenRM card level, not in the Super model card --- a multi-hop license-hygiene step that only becomes legible after joining the GenRM card with the upstream preference dataset.

    \item \textbf{Nemotron-Cascade SWE-SFT excludes repositories that appear in SWE-Bench Verified.}
    Where a typical model card simply claims ``decontamination was applied,'' the Nemotron-Cascade dataset card records the operational definition: SWE SFT instances whose source repository appears in \texttt{princeton-nlp/SWE-bench\_Verified} are dropped before SFT. Repository-level exclusion is stronger than text-level overlap removal and is good practice that papers often abstract away.

    \item \textbf{LLM360 MegaMath ships an exhaustive 11-benchmark n-gram-decontamination receipt.}
    \texttt{LLM360/MegaMath-Llama-3.2-1B} (and -3B) records 11 distinct \texttt{decontaminated\_against} edges to ASDiv, Hendrycks MATH, MathQA, MAWPS, MMLU-STEM, OCWCourses, SAT-Math, SVAMP, AIME 2024, AIME 2025, and AMC, each with a description naming the 12-gram filter applied to the question and answer fields. This per-benchmark, per-edge decontamination ledger is far richer than ``we decontaminated against standard benchmarks,'' and it makes future audits straightforward.

    \item \textbf{OLMo-3's Dolma 3 mix is published under ODC-BY with both complete and redacted variants.}
    Dolma 3 Mix is recorded as ``the primary mix licensed ODC-BY,'' with separate complete (non-redacted) and redacted variants of the academic-documents portion (the olmOCR science-PDF slice) so that downstream consumers can choose the appropriate license posture. This is another deliberate, multi-hop release-design choice surfaced by joining the model card, the dataset card, and the redaction notice.
\end{itemize}

\medskip

\noindent
Together, these findings show that recursive model--model dependency graphs are more than just a descriptive artifact; they provide a strong auditing layer and can surface key issues that were previously opaque due to the complexity of modern LLM development. Because \system{} uses reported information from public artifacts, these findings should be interpreted as a lower bound on the true dependency structure. Undocumented uses of private models, unreleased data mixtures, or internal filtering pipelines may introduce additional dependencies that are not recoverable from public evidence. An expanded set of supporting findings, including additional reward-model bootstrap chains, self-distillation patterns, snapshot-level dependencies, and per-stage judge multiplicities, is provided in the appendix.
